% Phase 2b work division — compiled with TeX Live (article + booktabs).
% Footer carries the AI-use acknowledgement, as required.
\documentclass[11pt,a4paper]{article}
\usepackage[T1]{fontenc}
\usepackage[utf8]{inputenc}
\usepackage[margin=22mm,top=20mm,bottom=24mm]{geometry}
\usepackage{booktabs}
\usepackage{microtype}
\usepackage{parskip}
\usepackage{fancyhdr}
\usepackage[colorlinks=true,linkcolor=black,urlcolor=black]{hyperref}
\pagestyle{fancy}\fancyhf{}
\renewcommand{\headrulewidth}{0pt}
\fancyfoot[L]{\scriptsize Team 23UG005 \textbullet\ Panel 4}
\fancyfoot[C]{\scriptsize Drafting assisted by an AI tool (Claude)}
\fancyfoot[R]{\scriptsize \thepage}

\begin{document}

\begin{center}
  {\Large\bfseries Phase 2b --- 100 UEs, three slices, KPI scorecard}\\[2mm]
  {\large Work division}\\[3mm]
  {\small Team 23UG005 \quad\textbullet\quad Panel 4}\\[1mm]
  {\small Anvita Arasavilli \quad Janhavi Nilesh Parate \quad Keerthi Devarajan Anuradha \quad Mukkara Rithvik Reddy}
\end{center}
\vspace{2mm}

\section*{What this phase delivers}

One hundred UEs are provisioned and attached in three classes --- 10 time-sensitive (URLLC),
20 bandwidth-hungry (eMBB) and 70 IoT devices (mMTC) that send ten bytes periodically and then
sleep. Each class becomes a network slice with its own KPI: latency for URLLC, throughput and its
share of system bandwidth for eMBB, and device density with delivery of small packets for mMTC.
Every KPI is measured on the running system and compared against published industry targets
(ITU-R M.2410 for IMT-2020 values, 3GPP TS 28.554 for the per-slice KPI definitions), weighted per
slice into a score, and exported to an Excel report that records every instance while one parameter
is varied. The orchestrator then grows from admitting demands to prioritising inside a slice and
deciding how resources are allocated. Fault diagnosis is deliberately left to the phase after this
one.

The work splits into four tracks that can run in parallel once the first one lands the slice and the
UEs. Names below are a proposal --- the tracks matter more than who holds them, and they can be
swapped as long as each track keeps one owner.

\section*{Track 1 --- Slice and scale (owner: Rithvik)}

Rithvik extends the running Open5GS deployment from two slices to three and from twenty UEs to a
hundred. This means an mMTC slice on the standard SST~3 with its own DNN, address pool, UPF and gNB
--- following the pattern already proven for the other two, where a shared gNB was what destroyed
isolation --- plus provisioning that hardcodes the 10/20/70 split, a UE supervisor that can start
and watch a hundred processes rather than twenty, and the deployment, gate and CI updated to the new
topology. The track's obligation to the others is a stack that reliably holds 100 attached UEs
across three slices, since every other track measures against it. Its own finding is the cost of
that scale: what a hundred simultaneous registrations do to the control plane, and where this laptop
stops.

\section*{Track 2 --- Traffic profiles and instruments (owner: Janhavi)}

Janhavi builds what each slice actually carries and how it is measured. URLLC needs small packets
sent at a fixed interval with a per-packet latency record, not an averaged ping; eMBB needs bulk
flows that saturate the slice so throughput is a real measurement rather than a configured number;
mMTC needs seventy devices each waking to send ten bytes and sleeping again, with every packet
accounted for so delivery can be computed. Seventy such devices amount to only a few kilobits per
second, which is exactly the point --- the mMTC instrument must count devices and delivered packets,
because measuring its bandwidth would say nothing. This track owns the traffic generators and the
exporters that turn their results into metrics, and it hands Track~3 a clean, named metric for every
KPI.

\section*{Track 3 --- KPI definitions, targets and scoring (owner: Keerthi)}

Keerthi decides what each KPI means, what number the industry expects, and how far the system is
from it. That is three pieces of work: writing the per-slice KPI set against the 3GPP catalogue so
the names and methods are defensible; recording the IMT-2020 target for each one with its source
clause, rather than a remembered figure; and defining the weights that turn several measurements
into one score per slice --- latency dominating URLLC, throughput dominating eMBB, density
dominating mMTC. The honest part of this track is the gap: URLLC's target is one millisecond and
this hardware delivers a few, so every score must carry the reason for its distance from the target.
A score without that explanation would mislead the reader, and this track owns making sure it never
does.

\section*{Track 4 --- Sweeps, reporting and review material (owner: Anvita)}

Anvita turns measurements into the report the review asks for. The requirement is an Excel record of
the KPIs at all instances while one parameter is varied, so this track designs the sweeps --- which
parameter moves (number of UEs, offered load, IoT reporting interval, packet size, or the split of
capacity between slices), which stay fixed, how long each point settles so the numbers are
comparable, and how a run is repeated. It then produces the workbook itself, with a sheet per sweep
and the per-slice scores alongside the raw KPIs, the Grafana panels that show the same story live,
and the slides and figures for the next review. Where a sweep shows something unexpected, this track
is responsible for saying so in the report rather than smoothing it.

\section*{Shared and later work}

The orchestrator extension --- priority within the URLLC slice, deciding how many slices a demand
set needs, and allocating from monitored load --- is shared between Track~1 and Track~3, because it
is an architecture decision carrying a measurement question, and it starts only once the KPI
definitions are settled. Fault diagnosis is the following phase and is not started here.

\section*{Sequence and hand-offs}

\begin{center}
\begin{tabular}{@{}clll@{}}
\toprule
\textbf{Order} & \textbf{What must happen} & \textbf{Owner} & \textbf{Blocks}\\
\midrule
1 & mMTC slice up and 100 UEs attached & Track 1 & every measurement\\
2 & Traffic profiles and exporters per class & Track 2 & KPI values\\
3 & KPI set, targets and weights agreed & Track 3 & scoring, sweeps\\
4 & Sweep design and workbook & Track 4 & the review report\\
5 & Orchestrator: priority and allocation & Tracks 1 + 3 & next review\\
\bottomrule
\end{tabular}
\end{center}

\vspace{2mm}
Tracks 2, 3 and 4 can begin on the current two-slice, twenty-UE stack before Track~1 finishes; none
of them needs to wait idle. The one rule that keeps the tracks from colliding is that a KPI is
defined in one place only --- the definitions file Track~3 owns --- and everyone else reads it.

\end{document}
